\scanpage{39}
\mspage{31}
% L1–6
Donc la cat. des triplets $(C,\Delta,b)$ d'un complexe cubique $C$, muni d'un bitriangle $\Delta$ et d'une base $b$ orthogonale à \unsure{celle-ci|celui-ci}, \uline{équivaut} à celle des $\Delta$ (car $V^*$ se réconstitue comme
\[
V^*\simeq J\amalg\{\mathrm{pt}\}
\]
% L7–9
Donc le groupe des autom. d'une telle structure est $\simeq\operatorname{Aut}(\Delta)\simeq\unsure{\mathfrak S|\mathfrak G}_2\cdot(\mathfrak S_3\times\mathfrak S_3)$ de card. 72.

% L10–19
Pour $C$ fixé, l'ens des couples $(\Delta,b)$ est de cardinal $\inserted{120\cdot6=}720$, donc pour $b\inserted{=\unsure{T'|F'|\mathbb F'}}$ fixé, l'ens des $\Delta$ correspondants est de cardinal $10=\binom63\frac12={}$\inserted{cardinal de l'}ens des partitions \inserted{$\{T'_1,T'_2\}$} de $T'$ en deux parties à 3 éléments. Or, pour une telle partition, on déduit un bitriangle orthogonal à $b$, en prenant \inserted{l'ens} des $(\{ij\})_{\substack{i\in T'_1\\j\in T'_2}}$, correspondant à l'ens \inserted{$T$} des 6 triangles \inserted{$t_{\sigma}={}$}$(\{i,\sigma i\})_{i\in T'_1}$, où $\sigma\in\operatorname{Isom}(T'_1,T'_2)$.
% L20–23. Two pen strokes in the left margin beside this passage.
\deleted{Ainsi,} En $=$ temps, on tombe sur une \uline{autodualité} dans la catégorie des ens de cardinal 6 munis d'une partition de type $(3,3)$ \ldots
% archive sheet number: 38
